

\section{Second Inverse Decomposition}
\label{sec: second inverse decomp}

In this section we refine the First Inverse Decomposition from \Cref{prop: second version}.
This version is best suited for establishing \Cref{Conj A} in the following section.

\subsection{The $\mathbf{v}$-operator}


\begin{definition}\rm \label{def: v negrita}
    Let $\lambda \in X^+$ and  $(j,h)$ be a pair of  integers such that $2\leq h \leq j$.
    Suppose that $\pint{\lambda}{\alpha_t}=0$ for all $j-(h-2)\leq t \leq j$.
    Let $\varkappa$ be the smallest non-negative integer such that $Q^{-\varkappa}_{j+1,h}(v^{h-1}_{j,h}(\lambda)) \in X^+$.
    \begin{itemize}
        \item If such an integer exists then we define $\mathbf{v}_{j,h}(\lambda) =  Q^{-\varkappa}_{j+1,h}(v^{h-1}_{j,h}(\lambda))$.
        \item Otherwise, we say that $\mathbf{v}_{j,h}(\lambda)$ is not defined. 
    \end{itemize}
     In the first case we call $\varkappa =\varkappa_{j,h}(\lambda) $ the integer associated to $\mathbf{v}_{j,h}(\lambda)$. 
\end{definition}


\begin{remark}\rm \label{rem: varkappa =0}
    We stress that if $v^{h-1}_{j,h}(\lambda)\in X^+$  then $\varkappa=0$ and $\mathbf{v}_{j,h}(\lambda) = v^{h-1}_{j,h}(\lambda)$.
    On the other hand, if $v^{h-1}_{j,h}(\lambda)$ is not defined, then $\mathbf{v}_{j,h}(\lambda)$ is automatically not defined.
\end{remark}

The following lemma asserts that, under repeated application of  $Q$-operators, the weight $v^{h-1}_{j,h}(\lambda)$ either transforms into the dominant weight $\mathbf{v}_{j,h}(\lambda)$, yielding an identity between the corresponding $\bfM$-elements, or else $v^{h-1}_{j,h}(\lambda)$ fails to reach the dominant chamber, in which case the associated $\bfM$-element vanishes.

\begin{lemma}\rm\label{lem: completar Q en caso h-1}
    Let  $\alpha_{i,j}\in \Phi^{\geq 2}$ and $h=j-i+1$.
    Let $\lambda\in X^{+}$ satisfying $\langle \lambda , \alpha_r \rangle  =0$, for all $i+1\leq r \leq j$.
    Suppose that $v^{h-1}_{j,h}(\lambda)$ is defined and set $Y= \Phi^{\succ \alpha_{i,j}} \cup \{ \alpha_{i-h+1,i}  \} $.
    Then, we have
\begin{equation}
    \bfM_{v^{h-1}_{j,h}(\lambda)}^{Y} =
    \begin{cases}
        q^{\varkappa}\bfM_{\mathbf{v}_{j,h}(\lambda)}^{Y}, & \mbox{if } \mathbf{v}_{j,h}(\lambda) \mbox{ is defined;} \\
        0, & \mbox{otherwise.}
    \end{cases}
\end{equation}
\end{lemma}




\begin{proof}
We first assume that $\mathbf{v}_{j,h}(\lambda)$ is defined.
If $\varkappa=0$ then  there is nothing to prove.
So that  we can assume $\varkappa>0$.
For $0 \leq k \leq \varkappa$ we define $\mu_k = Q^{-k}_{j+1,h}(v^{h-1}_{j,h}(\lambda))$. 
By combining \Cref{lem: desigualdad de varthetas} and \Cref{lem: thetas para el caso h-1}  we obtain for all $1\leq k \leq \varkappa$ that 
\begin{equation} \label{eq:S6A}
    \mu_k = \mu_{k-1} - \alpha_{j+k+1,j+k+h-1}, 
\end{equation}
or in  words, the $Q$-operator needed to pass from $\mu_{k-1} $ to $\mu_k$ consists of a unique root of height $h-1$. 
On the other hand,  \Cref{lem: thetas para el caso h-1} also implies that
\begin{equation}\label{eq:S6B}
  \mu_0= v^{h-1}_{j,h}(\lambda) = P^{h-1}_{j,h}(\lambda) - \sum_{r=1}^{h-1}\alpha_{i+r+1,j+r} =  P^{h-1}_{j,h}(\lambda) +\varpi_{i+1} -2\varpi_{j+1} +\varpi_{j+h}.
\end{equation}

Furthermore, by definition of $P^{h-1}_{j,h}(\lambda)$  and the assumptions on $\lambda$ we have 
\begin{equation}\label{eq:S6C}
    \pint{P_{j,h}^{h-1}(\lambda)}{\alpha_t} = \begin{cases}
    -1, & \mbox{if } t=i+1;\\
    0, & \mbox{if } t \in [i+2,j];\\
    \pint{\lambda}{\alpha_{j+1}} +1, & \mbox{if }  t=j+1.\\
        \pint{\lambda}{\alpha_t}, & \mbox{if } t\in [j+2,n];\\ 
    \end{cases}
\end{equation}                                                                                                                                                                                                                                                                                                             

Thus, combining \eqref{eq:S6A}, \eqref{eq:S6B}, and \eqref{eq:S6C} with \Cref{remark producto interno}, we obtain, for all 
$0\leq k<\varkappa$, that

\begin{equation}\label{eqS6 mu versus alpha}
\pint{\mu_k}{\alpha_t} = \begin{cases}
    0, & \mbox{if } t\in [i+1, j+k];\\
    -1,& \mbox{if } t=j+k+1.
\end{cases}
\end{equation}



Hence, for $0\leq k \leq \varkappa - 1$, we can apply \Cref{lem: free Q2} to $K=\Phi^{\succ\alpha_{i,j} } $, $\alpha_{i,j}$, $T=T_{k}= [i+1,j+k+1]$, $\mu_{k}$, $p' = j+k+1$ and $p$ be the unique integer such that $i+1\leq p\leq j$ and $p\equiv p' \mod{h-1}$ to get
    \begin{equation}
        \bfM_{\mu_{k}}^{Y} = q\bfM_{\mu_{k+1}}^{Y}.
    \end{equation}
    
Therefore,  $\bfM_{v^{h-1}_{j,h}(\lambda)}^{Y} =  \bfM_{\mu_{0}}^{Y} =q^{\varkappa}  \bfM_{\mu_{\varkappa}}^{Y} =  q^{\varkappa}\bfM_{\mathbf{v}_{j,h}(\lambda)}^{Y}$, as we wanted to show.

We now assume that $\mathbf{v}_{j,h}(\lambda)$ is not defined. 
Although this weight is not defined, the weight $\mu_{n-j-h+1}  $ does exist. 
Furthermore, it can be computed using \cref{eq:S6A} and \cref{eqS6 mu versus alpha} still holds. 
Thus, arguing as in the previous case we obtain
\begin{equation}\label{eq:S6D}
    \bfM_{v^{h-1}_{j,h}(\lambda)}^Y = q^{n-j-h+1} \bfM_{\mu_{n-j-h+1}}^{Y}.
\end{equation}

 On the other hand, by  \Cref{lem: free Q2} applied to $K=\Phi^{\succ \alpha_{i,j} } $, $\alpha_{i,j}$, $T=[i+1,n-h+2]$, $\mu_{n-j-h+1}\in X$, $p'=n-h+2$, and $p$ be the unique positive integer such that $i+1\leq p\leq j$ and $p\equiv p'\mod{h-1}$, we conclude that
    \begin{equation}\label{eq: final completar Q caso h-1}
        \bfM_{\mu_{n-j-h+1}}^{Y} =0.
    \end{equation}
 Therefore, the result follows by combining  \cref{eq:S6D} and  \eqref{eq: final completar Q caso h-1}. 
\end{proof}


%%%%%%%%%%%%%%%%%%%%%%%%%%%%%%%%%%%%%%%%%%%%%%%%%%%%%%%%%%%%
%%%%%%%%%%%%%%%%%%%%%%%%%%%%%%%%%%%%%%%%%%%%%%%%%%%%%%%%%%%%
\subsection{The $\mathbf{P}$-operator}
%%%%%%%%%%%%%%%%%%%%%%%%%%%%%%%%%%%%%%%%%%%%%%%%%%%%%%%%%%%%
%%%%%%%%%%%%%%%%%%%%%%%%%%%%%%%%%%%%%%%%%%%%%%%%%%%%%%%%%%%%


\begin{definition}\rm  \label{def: P negrita}
  Let $(h,k)$ be a pair of integers such that $2\leq h \leq k \leq n$.
  Let $\lambda \in X_{k}^+(L)$. 
  Let $x$ be the smallest positive integer such that $P^{x}_{k,h}(\lambda)\in X_{k}^+(L)$.
  If such an integer exists, then we define $\mathbf{P}_{k,h}(\lambda) =  P^{x}_{k,h}(\lambda)$. 
  Otherwise, we say that  $\bfP_{k,h}(\lambda)$ is not defined. 
\end{definition}

The main goal of this section is to prove the next lemma. 

\begin{lemma}\rm\label{lem: bfP}
  Let $\alpha_{i,j}\in \Phi^{\geq 2}$ and $h=j-i+1$.
  Let  $\lambda \in X_{j}^+(L)$ and  $\mu = P^{h}_{j,h}(\lambda)$.
  If $\pint{\mu}{\alpha_{i}} = -1$, then
    \begin{equation}\label{eq: bfP}
        \bfM_{\mu}^{\succ \alpha_{i,j}} = \left\{ \begin{array}{ll}
            q^{D(\mathbf{P}_{i-1,h})(\mu)}\bfM_{\mathbf{P}_{i-1,h}(\mu)}^{\succ \alpha_{i,j}}, & \text{if } \mathbf{P}_{i-1,h}(\mu) \text{ is defined;} \\[2pt]
            0, & \text{otherwise.}
        \end{array}\right.
    \end{equation}
\end{lemma}

Before embarking in the  proof of \Cref{lem: bfP} we need some preparatory results. 

\begin{lemma}\rm\label{lem: blackP caso 0}
    Let $\alpha_{i,j}\in \Phi^{\geq 2}$  and  $h=j-i+1$.
    Suppose that  $i+1\equiv 0 \mod h$.  
    Let $b$ be an integer such that $h \leq b \leq i$ and $ b\equiv 0 \mod{h}$.
    Let $\mu\in X_{h-1}^{+}(L) $  be such that $\pint{\mu}{\alpha_t}=-\delta_{t,b}$ for $t\in [h,i]$. 
    Then we have
    \begin{equation}\label{eq: first caso 0 blackP}
        \bfM_{\mu - \alpha_{i,j}}^{\succ \alpha_{i,j}} = 0.
    \end{equation}
\end{lemma}


\begin{proof}
We proceed by induction on $m = \pint{\mu}{\alpha_{h-1}}$.
If $m=0$ then $\pint{\mu-\alpha_{i,j}}{\alpha_t}=0$ for all $1\leq t \leq i-h$ and $t\equiv i-h \equiv -1 \mod h$. 
By definition of the $P_{i-1,h}$-operator, it follows that  $P_{i-1,h}(\mu-\alpha_{i,j})$ is not defined.
Then,  by \Cref{lem: free P} applied to $K =\Phi^{\succ \alpha_{i,j}}$, $\alpha_{i,j}$, $p = i = p'$ and the weight $\mu- \alpha_{i,j}$ we obtain $\bfM_{\mu-\alpha_{i,j}}^{\succ \alpha_{i,j}} = 0$.

We now suppose that $m\geq 1$ and assume that \eqref{eq: first caso 0 blackP} holds for $m-1$.
In this case $P_{i-1,h}(\mu-\alpha_{i,j})$ is defined.
In fact, we have
\begin{equation}
    \mu'\coloneqq P_{i-1,h}(\mu-\alpha_{i,j}) = \mu-\alpha_{i,j} -\alpha_{h-1,i-1}=\mu-\alpha_{h-1,j}= P_{j,h}(\mu) .
\end{equation}
Then, applying \Cref{lem: free P} with the same parameters as before, we get 
\begin{equation}\label{eq: first caso 0 blackP proof 1}
    \bfM_{\mu-\alpha_{i,j}}^{\succ \alpha_{i,j}} = q^{D(P_{i-1,h})(\mu  - \alpha_{i,j})}\bfM_{\mu'}^{\succ\alpha_{i,j}}=  q^{D(P_{j,h})(\mu  )-1}\bfM_{\mu'}^{\succ\alpha_{i,j}}.
\end{equation}

On the other hand, by the conditions imposed on $b$ and $\mu$ it is easy to see that $P_{b-1,h}(\mu')$ is not defined.
 Therefore, another application of  \Cref{lem: free P}, but now for the parameters
\begin{equation}
\displaystyle    K =\Phi^{\succ\alpha_{i-1,j-1}}, \quad \alpha_{i-1,j-1}, \qquad p = i-h+1, \qquad p'= b, \qquad k = \frac{i-b+1}{h}-1 \quad \mbox{ and } \quad \mu',
\end{equation}
 yields $\bfM_{\mu'}^{\succ \alpha_{i-1,j-1}} =0$.
 By the definition of $\bfM$-elements  we have
 \begin{equation} \label{eq: first caso 0 blackP proof 1 AA}
     0 = \bfM_{\mu'}^{\succ \alpha_{i-1,j-1}}  = \bfM_{\mu'}^{\succ \alpha_{i,j}} -q \bfM_{\mu' -\alpha_{i,j}}^{\succ \alpha_{i,j}}.
 \end{equation}
Thus, by combining \eqref{eq: first caso 0 blackP proof 1} and \eqref{eq: first caso 0 blackP proof 1 AA}, we obtain


\begin{equation}\label{eq: first caso 0 blackP proof 1 AAA}
    \bfM_{\mu-\alpha_{i,j}}^{\succ \alpha_{i,j}} = q^{D(P_{j,h})(\mu)}\bfM_{\mu' -\alpha_{i,j}}^{\succ\alpha_{i,j}}.
\end{equation}

Finally, since $\mu'$ satisfies the hypothesis of the lemma and $\pint{\mu'}{\alpha_{h-1}} =m-1$ we can apply our inductive hypothesis to conclude that the right-hand side of \cref{eq: first caso 0 blackP proof 1 AAA} vanishes, which completes the proof. 
\end{proof}








% \begin{lemma}\rm \label{lem: 0s de PPPPPPP}
%     Let  $\alpha_{i,j}\in \Phi^{\geq 2}$ and $h=j-i+1$.
%     Let $\lambda\in X_{j}^+(L)$ and for $x\geq 0$ we set  $\lambda_{x} = P^x_{j,h}(\lambda)$.
%     Let $k\geq h+1$ be an integer such that  $\lambda_{k}$ is defined.
%     Suppose that $\pint{\lambda_{x}}{\alpha_{j-(x-1)}} = -1$ for all $h\leq x<k$.
%     Then, we have
%     \begin{equation}\label{eq: 0s de PPPPPPP}
%         \pint{\lambda_{k-1}}{\alpha_{t}} =0,\quad\text{for all }\quad t\in [\varrho_{k}+h,i]\setminus \{j-(k-2)\},
%     \end{equation}
%    where $\varrho_{k} = \varrho_{j-(k-1),h}(\lambda_{k-1})$.
   
%   Additionally, if $\varrho_{k}+h \leq j-k+1 \leq i$, then $\pint{\lambda_{k}}{\alpha_{j-k+1}} = -1$.
% \end{lemma}



\begin{lemma}\rm \label{lem: 0s de PPPPPPP}
    Let  $\alpha_{i,j}\in \Phi^{\geq 2}$ and $h=j-i+1$.
    Let $\lambda\in X_{j}^+(L)$ and for $x\geq 0$ we set  $\lambda_{x} = P^x_{j,h}(\lambda)$.
    Let $k\geq h+1$ be an integer such that  $\lambda_{k}$ is defined and $\lambda_{u}\notin X_{j}^+(\lambda) $  for $h\leq u<k$.
    Then, we have
    \begin{equation}\label{eq: 0s de PPPPPPP}
        \pint{\lambda_{k-1}}{\alpha_{t}} =0,\quad\text{for all }\quad t\in [\varrho_{k}+h,i]\setminus \{j-(k-2)\},
    \end{equation}
   where $\varrho_{k} = \varrho_{j-(k-1),h}(\lambda_{k-1})$.
\end{lemma}



\begin{proof}
This result  follows by an inductive argument  on $k$ that  combines  \Cref{lem: desigualdad de varrhos},  \Cref{lem: ceros por Ps} and the definition of the $P$-operator. 
We left the details to the reader.
\end{proof}


% \begin{lemma}\rm\label{lem: first blackP}
%     Let $\alpha_{i,j}\in \Phi^{\geq 2}$ and $h=j-i+1$. 
%     Let  $b$ be an integer such that $h+1 \leq b \leq i$.
%     Let $\lambda\in X_j^+(L)$ and set  $\lambda_{x} = P^{x}_{j,h}(\lambda)$ for $x\geq 0$.
%     Suppose that $\lambda_{j-b+1}$ is defined and $\pint{\lambda_{x}}{\alpha_{j-x+1}}=-1$ for $h\leq x\leq j-b+1$.
    
%     Then
%     \begin{equation}\label{eq: first blackP}
%         \bfM_{\lambda_{j-b+1}}^{\succ \alpha_{i,j}} = \left\{\begin{array}{ll}
%             q^{\D{P,b-1,h}(\lambda_{j-b+1})}\bfM_{\lambda_{j-b+2}}^{\succ \alpha_{i,j}},&\text{if}\quad \lambda_{j-b+2}\quad\text{is defined};\\
%             0, &\text{otherwise}.
%         \end{array}\right. 
%     \end{equation}
% \end{lemma}


\begin{lemma}\rm\label{lem: first blackP} 
    Let $\alpha_{i,j}\in \Phi^{\geq 2}$ and $h=j-i+1$. 
    Let $\lambda\in X_j^+(L)$ and set  $\lambda_{x} = P^{x}_{j,h}(\lambda)$ for $x\geq 0$.
    Let  $u$ be an integer such that $h \leq u \leq i-1$.
    Suppose that $\lambda_{u}$ is defined and $\lambda_{x}\notin X_{j}^+(\lambda) $  for $h\leq x<u$. 
    Then, we have
    \begin{equation}\label{eq: first blackP}
        \bfM_{\lambda_{u}}^{\succ \alpha_{i,j}} = \left\{\begin{array}{ll}
            q^{\D{P,j-u,h}(\lambda_{u})}\bfM_{\lambda_{u+1}}^{\succ \alpha_{i,j}},&\text{if } \lambda_{u+1}\text{ is defined};\\
            0, &\text{otherwise}.
        \end{array}\right. 
    \end{equation}
\end{lemma}

\begin{proof}
Let $b= i-u+h$. 
We remark that $b-1=j-u$. 
We split the proof in four cases.

\medskip
\textbf{Case A. }   $\lambda_{u+1}$ is defined and $b\equiv i+1 \mod h$.
\medskip

We recall that $\lambda_{u+1}=P_{b-1,h}(\lambda_{u})$. 
Let $\varrho_{u+1} = \varrho_{b-1,h}(\lambda_{u})$ be the integer associated to $P_{b-1,h}(\lambda_{u})$.

By \Cref{lem: 0s de PPPPPPP}  applied to $k = u+1$ we obtain
 $\pint{\lambda_{u}}{\alpha_{t}} = 0$ for all 
 $t\in T\coloneq[\varrho_{u+1}+h,i]\setminus \{b\}$.
    Furthermore, since $\lambda_{u}\notin X_{j}^+(L)$ we have
   $\pint{\lambda_{u}}{\alpha_{b}} = -1$.

 We  proceed by induction on $D = \D{P,b-1,h}(\lambda_{u})$.

Suppose that $D=1$.
In this case we have $\varrho_{u+1}=i-u=b-h$ and $T=[b+1,i] $.  
Therefore, we can apply  \Cref{lem: left solving} to $\alpha_{i,j}$,  $K=\Phi^{\succ \alpha_{i,j}}$, $p=i-h+1$,  $p^{\prime} = b$ and $\lambda_{u}$ to get 
   \begin{equation}
      \bfM_{\lambda_{u}}^{\succ \alpha_{i,j}}= q\bfM_{\lambda_{u} -\alpha_{b-h,b-1}}^{\succ \alpha_{i,j}} = q\bfM_{\lambda_{u+1}}^{\succ \alpha_{i,j}},  
   \end{equation}
   which proves the lemma for $D=1$.    
   
 We now suppose that  $D>1 $ and that the lemma holds for $D-1$. 
Let $a \coloneqq  \varrho_{u+1} +h$. 
We notice that $a\equiv i-u \mod h$ and therefore $a\equiv b \mod h$. 
Furthermore, since $D>1$ we have $a<b$ and since $b\neq i$ we have $b<i$. 
Thus,  we can apply \Cref{lem: primer paso PPPPPP} to $\alpha_{i,j}$, $a$ and $b$ to obtain $\bfM_{\lambda_{u}}^{Y}=0$, where $Y=\Phi^{\succ \alpha_{i,j}} \cup \{ \alpha_{a-h,a-1}  \}$.
By definition of $\bfM$-elements we have
\begin{equation}
    \bfM_{\lambda_{u}}^{\succ \alpha_{i,j}} = q\bfM_{\lambda_{u} - \alpha_{a-h,a-1}}^{\succ \alpha_{i,j}}.
\end{equation}
On the other hand, we recall that 
\begin{equation}
    \lambda_{u+1} =P_{b-1,h}(\lambda_{u})=\lambda_{u}-\alpha_{\varrho_{u+1}
    ,b-1} =\lambda_{u}-\alpha_{a-h,b-1} =(\lambda_{u}-\alpha_{a-h,a-1})-\alpha_{a,b-1}.
\end{equation}
Since $ \pint{\lambda_{u}-\alpha_{a-h,a-1}}{\alpha_{a}}= 1$ we conclude that 
$ P_{b-1,h}(\lambda_{u}-\alpha_{a-h,a-1}) =  \lambda_{u+1}$.
Furthermore, we have $D(P_{b-1,h})(\lambda_{u} - \alpha_{a-h,a-1}) = D-1$. 
Therefore, our inductive hypothesis gives
   \begin{equation}
       \bfM_{\lambda_{u}}^{\succ \alpha_{i,j}} = q\bfM_{\lambda_{u} - \alpha_{a-h,a-1}}^{\succ \alpha_{i,j}} = qq^{D-1}\bfM_{\lambda_{u+1}}^{\succ \alpha_{i,j}}= q^{D}\bfM_{\lambda_{u+1}}^{\succ \alpha_{i,j}}.
   \end{equation}




\medskip
\textbf{Case B. }  $\lambda_{u+1}$ is not defined and $b\equiv i+1 \mod h$.
\medskip


Since $\lambda_{u}$ is defined and  $\lambda_{u+1}$ is not, \Cref{lem: desigualdad de varrhos} implies  $\varrho_{u} = \varrho_{b,h}(\lambda_{u-1}) = 1$. 
In particular, we have $b\equiv 0 \mod h$.
In addition, by definition of the $P$-operator we have $\pint{\lambda_{u}}{\alpha_{t}}=0$ for all $1\leq t \leq b-h $ with $t\equiv b-h \equiv 0 \mod h$. 
In particular, $\pint{\lambda_{u}}{\alpha_{h}}=0$. 

On the other hand, by  \Cref{lem: 0s de PPPPPPP} applied to $k=u$ we obtain 
$\pint{\lambda_{u-1}}{\alpha_{t}}$ for all $t\in [h+1, i] \setminus \{b+1\}$. 
Furthermore, $\lambda_{u-1}\notin X_j^{+}(L)$ implies $\pint{\lambda_{u-1}}{\alpha_{b+1}}=-1$. 
Since $\lambda_u= \lambda_{u-1}- \alpha_{1,b}$ we obtain  $\pint{\mu_u}{\alpha_t}=0$ for all $t\in  [h+1,i]\setminus\{b\}$ and $\pint{\lambda_{u}}{\alpha_{b}} = -1$.

Summing up, we obtain that $\pint{\lambda_u}{\alpha_t}=0$ for all $t\in  [h,i]\setminus\{b\}$ and $\pint{\lambda_{u}}{\alpha_{b}} = -1$.
Then, we can apply  \Cref{lem: free P} to $\alpha_{i-1,j-1}$, $K=\Phi^{\succ \alpha_{i-1,j-1}}$, $p =i+1-h$, $p'=b$ and $\lambda_{u}$ to obtain $\bfM_{\lambda_{u}}^{\succ \alpha_{i-1,j-1}} = 0$. 
We remark that our use of \Cref{lem: free P} is justified since we use $\alpha_{i-1,j-1}$ rather than $\alpha_{i,j}$, so that our $p$ belongs to the admissible range.
By decomposing $\bfM_{\lambda_{u}}^{\succ \alpha_{i-1,j-1}} = 0$ we get
    \begin{equation}\label{eq: blackP proof 1}
        \bfM_{\lambda_{u}}^{\succ \alpha_{i,j}} =q \bfM_{\lambda_{u} - \alpha_{i,j}}^{\succ \alpha_{i,j}}.
    \end{equation}

Since $b\equiv 0 \mod h$, \Cref{lem: blackP caso 0}  applied to $\mu = \lambda_{u}$ yields  $\bfM_{\lambda_{u} - \alpha_{i,j}}^{\succ \alpha_{i,j}} = 0$. 
Therefore, the result follows by replacing this in \eqref{eq: blackP proof 1}.

\medskip
\textbf{Case C.}  $\lambda_{u+1}$ is defined and $b \not \equiv i+1 \mod h$.
\medskip

As in Case A we obtain  $\pint{\lambda_{u}}{\alpha_{t}} = 0$ for all $t\in [\varrho_{u+1}+h,i]\setminus \{b\}$ and  $\pint{\lambda_{u}}{\alpha_{b}} = -1$.
    Let $p$ be  the unique integer such that  $i-h+1< p\leq i$ and $p\equiv b \mod h$.
    Let $d$ be the integer defined by $\varrho_{u+1} = p-(d+1)h$ and define
    \begin{equation}
        T = \bigcup_{m = 0}^{d}[p-mh,i-mh]\subset [\varrho_{u+1} +h,i].
    \end{equation}
   We apply  \Cref{lem: free P} to  $K=\Phi^{\succ \alpha_{i,j}}$, $\alpha_{i,j}$, $T$, $\lambda_{j-b+1}$, $p' = b$ and $p$ to obtain
\begin{equation}
    \bfM_{\lambda_u}^{\succ\alpha_{i,j}} =q^{D(P_{b-1,h})(\lambda_{u})}  \bfM_{P_{b-1,h}(\lambda_{j-b+1})}^{\succ \alpha_{i,j}} =q^{D(P_{b-1,h})(\lambda_{u})}  \bfM_{\lambda_{u+1}}^{\succ \alpha_{i,j}} ,
\end{equation}
as we wanted to show. 

\medskip
\textbf{Case D.} $\lambda_{u+1}$ is not defined and  $b \not \equiv i+1 \mod h$.
\medskip

As in Case B we get $\varrho_{u}=1$,    $\pint{\lambda_u}{\alpha_t}=0$ for all $t\in  [h+1,i]\setminus\{b\}$ and $\pint{\lambda_{u}}{\alpha_{b}} = -1$.   
Let $p$ be the unique integer such that  $i-h+1\leq p\leq i$ and $p\equiv b \mod h$. 
Then,  by \Cref{lem: free P} applied to $\alpha_{i,j}$, $K=\Phi^{\succ \alpha_{i,j}}$,  $\lambda_{u}$, $p' = b$ and $p$ we obtain that $\bfM_{\lambda_{u}}^{\succ\alpha_{i,j}}=0$.
\end{proof}


We are now in position of proving the main result in this section. 



\begin{myproof}[Proof of \Cref{lem: bfP}]
For $x \geq 0$, set $\lambda_x = P^x_{j,h}(\lambda)$.
We stress that $\mu=\lambda_h $.

We first consider the case when $\mathbf{P}_{i-1,h}(\mu)$ is defined.
By hypothesis we have $\lambda_{h} \notin X_{j}^+(L)$.
Let $s\geq h+1$ be the integer such that $\lambda_s = \mathbf{P}_{i-1,h}(\mu) $.
Let $h+1\leq u <s$. 
By the minimality of $s$ we have $\lambda_{u} \notin X_{j}^+(L)$. 


Then, by  \Cref{lem: first blackP} we have
\begin{equation}
    \bfM_{\lambda_{u}}^{\succ \alpha_{i,j}}
    = q^{D(P_{j-u,h})(\lambda_{u})}
      \bfM_{\lambda_{u+1}}^{\succ \alpha_{i,j}}
\end{equation}
for all $h \leq u <s$. It follows that
\begin{equation}
 \displaystyle   \bfM_{\mu}^{\succ \alpha_{i,j}} =   \bfM_{\lambda_h}^{\succ \alpha_{i,j}}
    = q^{D}
      \bfM_{\lambda_s}^{\succ \alpha_{i,j}}
    = q^{D(\mathbf{P}_{i - 1,h})(\mu)}
      \bfM_{\mathbf{P}_{i - 1,h}(\mu)}^{\succ \alpha_{i,j}},
\end{equation}
where $\displaystyle D=\sum_{u =h}^{s-1} D(P_{j-u,h})(\lambda_{u})$.
This proves the lemma when $\mathbf{P}_{i - 1,h}(\mu)$ is defined.


We now address the case when $\mathbf{P}_{i-1,h}(\mu)$ is not defined.
Let $s\geq h+1$ be the smallest  integer such that $\lambda_s$ is not defined.
Arguing as in the previous paragraph we obtain
\begin{equation}
    \bfM_{\mu}^{\succ \alpha_{i,j}} =   \bfM_{\lambda_h}^{\succ \alpha_{i,j}} = q^E \bfM_{\lambda_{s-1}}^{\succ \alpha_{i,j}}, 
\end{equation}
for some integer $E$. 
Then another application of \Cref{lem: first blackP} yields $\bfM_{\lambda_{s-1}}^{\succ \alpha_{i,j}}=0$, which proves the lemma in this case. 
\end{myproof}


%%%%%%%%%%%%%%%%%%%%%%%%%%%%%%%%%%%%%%%%%%%%%%%%%%%%%%%%%%%%%%%%%%%%
\subsection{Mixed operators}
%%%%%%%%%%%%%%%%%%%%%%%%%%%%%%%%%%%%%%%%%%%%%%%%%%%%%%%%%%%%%%%%%%%%

In this section we establish a relation between the weights 
$P^{h-1}_{j,h}(\lambda)$, 
$\mathbf{v}_{j,h}(\lambda)$,
$P^{h}_{j,h}(\lambda)$ and
$P_{i,h}\mathbf{v}_{j,h}(\lambda)$
(when all or some of them exist).
This is the key to pass from the First Inverse Decomposition to the Second Inverse Decomposition. 

\begin{lemma}\rm\label{lem: ultimo P}
    Let $\lambda\in X^{+}$,  $\alpha_{i,j}\in \Phi^{\geq 2}$ and $h=j-i+1$. 
    Suppose that $\langle \lambda , \alpha_r \rangle  =0$, for all $i\leq r \leq j$.
    and let 
\begin{equation}
    \begin{array}{lll}
      d_{1} =D(P^{h-1}_{j,h})(\lambda)+1   & \mbox{ }\qquad \mbox{ }\qquad    & d_{2}=D(\mathbf{v}_{j,h})(\lambda)+1 \\
      d_{3} = D(P^{h}_{j,h})(\lambda)   & \mbox{ }\qquad  \mbox{ }\qquad    &  d_{4}=D(P_{i,h}\mathbf{v}_{j,h})(\lambda)
    \end{array}
\end{equation}
    \begin{enumerate}[(a)]
        \item Suppose that $P^{h}_{j,h}(\lambda)$ and $\mathbf{v}_{j,h}(\lambda)$ are defined, then\label{item P existe}
    \begin{equation} \label{eq: ultimo P existe}
      q^{d_{1}} \bfM_{P^{h-1}_{j,h}(\lambda) - \alpha_{i-h+1,i}}^{\succ \alpha_{i,j}} - q^{d_{2}} \bfM_{\mathbf{v}_{j,h}(\lambda)-\alpha_{i-h+1,i}}^{\succ \alpha_{i,j}} = q^{d_{3}} \bfM_{P^{h}_{j,h}(\lambda)}^{\succ \alpha_{i,j}} - q^{d_{4}} \bfM_{P_{i,h}\mathbf{v}_{j,h}(\lambda)}^{\succ \alpha_{i,j}}.
    \end{equation}
        \item Suppose that $\mathbf{v}_{j,h}(\lambda)$ is defined, but $P^{h}_{j,h}(\lambda)$ is not, then\label{item P no existe}
    \begin{equation}\label{eq: ultimo P no existe}
        q^{d_{1}} \bfM_{P^{h-1}_{j,h}(\lambda) - \alpha_{i-h+1,i}}^{\succ \alpha_{i,j}} - q^{d_{2}} \bfM_{\mathbf{v}_{j,h}(\lambda)-\alpha_{i-h+1,i}}^{\succ \alpha_{i,j}} = 0.
    \end{equation}
        \item Suppose that $P^{h}_{j,h}(\lambda)$ is defined, but $\mathbf{v}_{j,h}(\lambda)$ is not, then\label{item V no existe}
    \begin{equation}\label{eq: ultimo P  y v no existe}
        \bfM_{P^{h-1}_{j,h}(\lambda) - \alpha_{i-h+1,i}}^{\succ \alpha_{i,j}} = q^{D(P_{i,h})(P_{j,h}^{h-1}(\lambda))} \bfM_{P^{h}_{j,h}(\lambda)}^{\succ \alpha_{i,j}}.
    \end{equation}
        \item Suppose that $P^{h}_{j,h}(\lambda)$ and $\mathbf{v}_{j,h}(\lambda)$are not defined, then\label{item p y v no existen}
    \begin{equation}\label{eq: ultimo P, ambos no existen}
        \bfM_{P^{h-1}_{j,h}(\lambda) - \alpha_{i-h+1,i}}^{\succ \alpha_{i,j}} = 0.
    \end{equation}
    \end{enumerate}
\end{lemma}

\begin{proof}
    We only prove \cref{eq: ultimo P existe}, the other three cases being analogous. 

We  notice that $P_{j,h}^{h}(\lambda) = P_{i,h}(P_{j,h}^{h-1} (\lambda) ) $. 
Let  $\varrho=\varrho_{i,h}(P_{j,h}^{h-1} (\lambda)) $. 
By definition if $d=D(P_{i,h} )(P_{j,h}^{h-1} (\lambda))$ then $\varrho = i -dh+1$. 
Furthermore, we have the 
\begin{equation}\label{eq:degreesA}
    d_1+d=d_3+1.
\end{equation}
We also notice that $\pint{P_{j,h}^{h-1} (\lambda)}{\alpha_t} =     \pint{\mathbf{v}_{j,h}(\lambda)}{\alpha_t}  $  for all $1\leq t \leq i$. 
It follows that $\varrho = \varrho_{i,h} (  \mathbf{v}_{j,h}(\lambda) )$.

On the other hand, let $\varkappa = \varkappa_{j,h}(\lambda) $.
Then, by definition of $\mathbf{v}_{j,h}(\lambda)$ we have 
\begin{equation}\label{eq: relation P con v negrita}
    \mathbf{v}_{j,h}(\lambda) = P^{h-1}_{j,h}(\lambda) -\sum_{u=1}^{h-1+\varkappa}  \alpha_{i+1+u,i+1+u+(h-2)}.
\end{equation}
It follows that 
\begin{equation}\label{eq:degreesB}
    d_2=d_1 + \varkappa +h -1  \qquad \mbox{and} \qquad d_4=d_3 + \varkappa +h -1.
\end{equation}


We prove the following.

\begin{claim}\label{claim complicado}
For all  $\varrho < a \leq i-h+1$ with $a\equiv i+1 \mod h$ we have
\begin{equation}
   \bfM_{P_{j,h}^{h-1}(\lambda) - \alpha_{a,i}}^{\succ \alpha_{i,j}} - q^{\varkappa+h-1} \bfM_{\mathbf{v}_{j,h}(\lambda)-\alpha_{a,i}}^{\succ \alpha_{i,j}}  = q  \left( \bfM_{P_{j,h}^{h-1}(\lambda)-\alpha_{a-h,i}}^{\succ \alpha_{i,j}} - q^{\varkappa+h-1} \bfM_{\mathbf{v}_{j,h}(\lambda)- \alpha_{a-h,i}}^{\succ \alpha_{i,j}} \right).
\end{equation}   
\end{claim}

\begin{proof}
    Let $\mu= P_{j,h}^{h-1}(\lambda) - \alpha_{a,i}$ and $\mu' = \mathbf{v}_{j,h}(\lambda)-\alpha_{a,i}$. 
By combining \Cref{lem: 0s de PPPPPPP}, the hypothesis on $\lambda$ in the lemma, the inequality $\varrho \leq a-h$ and the definition of $\varkappa$ we obtain 
\begin{equation}
    R_t(\mu) = \left\{\begin{array}{rl}
        1, & \text{if } t=a;\\
        0, & \text{if } a+1\leq t < i;\\
        1, & \text{if } t=i;\\
        0, & \text{if } i+1\leq t\leq j+\varkappa, \text{ and } t\neq j+1;\\
        -1, &\text{if } t = j+1\leq j+\varkappa.
    \end{array}\right.
    \qquad  R_t(\mu') = \left\{\begin{array}{rl}
        1, & \text{if } t=a;\\
        0, & \text{if } a+1\leq t < i;\\
        1, & \text{if } t=i;\\
        -1, & \text{if } t= i+1;\\
        0, & \text{if } i+2\leq t\leq j+\varkappa.
    \end{array}\right.
\end{equation}

    
    Let $Y=\Phi^{\succ\alpha_{i,j}} \cup \{ \alpha_{a-h,a-1} \}$ and $k\geq 0$ defined by the equality $i-h+1-a=kh$.  
    We notice that both $\mu$ and $\mu'$ satisfies the hypothesis of \Cref{lem: fourth string lemma} with $a$ and $k$ as above. Therefore,  using the notation in that lemma, we obtain 
    \begin{equation} \label{eq: alternativa 1}
        \bfM_{\mu}^{Y} = q^{\mathfrak{R}_k}\bfM_{\mu- \mu_k}^{Y_{a,k}}  \qquad \mbox{and} \qquad   \bfM_{\mu'}^{Y} = q^{\mathfrak{R}_k}\bfM_{\mu'- \mu_k}^{Y_{a,k}}.
    \end{equation}
    We stress that both the weight $\mu_k$ and the integer $\mathfrak{R_k}$ depends heavily on some of the components of $\mu$ and $\mu'$. 
    However, in our setting the components of $\mu$ and $\mu'$ involved in the computation of   $\mu_k$ and $\mathfrak{R}_k$ coincide. 
    More precisely, we have 
    \begin{equation}
        \mu_0 = \sum_{t=0}^{h-2} \alpha_{i-t,j-t}, \qquad \mu_{s} = \mu_{s-1} + \sum_{t=0}^{(s+1)h-2}    \alpha_{i-t,j-t}    \qquad \mbox{and} \qquad \mathfrak{R}_{k} =  \frac{k+1}{2}(i-a+h-1) .
    \end{equation}
In particular, we have
\begin{equation}
    \pint{\mu-\mu_k}{\alpha_{i-s}}= \pint{\mu'-\mu_k}{\alpha_{i-s}} = \left\{ \begin{array}{rl}
        -(k+2) , & \mbox{if } s=0;  \\
         0 ,& \mbox{if } 0<s\leq k. 
    \end{array}
    \right.
\end{equation}
    
Then, using  \eqref{eq: alternativa 1}  and \Cref{lem: weyl group over M elements} for the element $s_{i-k}\cdots s_{i-1}s_i $, we obtain 

        \begin{equation} \label{eq: alternativa 2}
        \bfM_{\mu}^{Y} = (-1)^{k+1}q^{\mathfrak{R}_k}\bfM_{\mu- \mu_k  +\gamma_k}^{Y_{a,k}}  \qquad \mbox{and} \qquad   \bfM_{\mu'}^{Y} = (-1)^{k+1}q^{\mathfrak{R}_k}\bfM_{\mu'- \mu_k +\gamma_k}^{Y_{a,k}}, 
    \end{equation}
where $\gamma_{k} =\displaystyle \sum_{s=0}^k (k+1-s) \alpha_{i-s}$. 

To simplify notation we set $ \Upsilon_k \coloneqq \mu- \mu_k  +\gamma_k$ and $\Upsilon_k'\coloneqq\mu'- \mu_k  +\gamma_k$. 
We have
\begin{equation} \label{eq: upsilon }
    \pint{\Upsilon_k}{\alpha_t} = \left\{  
    \begin{array}{rl}
      0,    & \mbox{if } i-k \leq t \leq i; \\
      -(k+2),    & \mbox{if } t=i+1; \\
      0,  & \mbox{if }  i+2\leq  t \leq j+\varkappa \mbox{ and } t\neq j+1;\\
      k+2, & \mbox{if }  t=j+1 \leq j+\varkappa.
    \end{array} \right.
\end{equation}
and 
\begin{equation} \label{eq: upsilon prime}
  \pint{\Upsilon_k'}{\alpha_t} = \left\{  
    \begin{array}{rl}
      0,    & \mbox{if } i-k \leq t \leq i; \\
      -(k+1),    & \mbox{if } t=i+1; \\
      0,  & \mbox{if }  i+2\leq  t \leq j+\varkappa \mbox{ and } t\neq j+1;\\
      k+1, & \mbox{if }  t=j+1 \leq  j+\varkappa.
    \end{array} \right.
\end{equation}


On the other hand, we observe that  since $b_k=i-1$ we have $Y_{a,k}=\Phi^{\succ \alpha_{a-h,a-1}} \setminus A_k$, where $A_k$ is the trapezoid at height $h$ with top-left corner $\alpha_{i-(h+k)+1,i+1}$ and right-bottom corner $\alpha_{i,j}$. 
We emphasize that $h' \coloneq \operatorname{ht} (\alpha_{i-(h+k)+1,i+1}) = h+k+1$. 
So that $h'-h=k+1 $. 

Therefore we are in position to apply \Cref{coro: push pyramid to the right}. 
We obtain
\begin{equation}
    \bfM_{\Upsilon_k}^{Y_{a,k}} =q^{k+2} \bfM_{\Upsilon_k-(k+2)\alpha_{i+2,i+2+(h-2)}}^{\Phi^{\succ \alpha_{a-h,a-1}} \setminus A_{k}(1)} \qquad \mbox{and} \qquad  \bfM_{\Upsilon_k'}^{Y_{a,k}} =q^{k+1} \bfM_{\Upsilon_k'-(k+1)\alpha_{i+2,i+2+(h-2)}}^{\Phi^{\succ \alpha_{a-h,a-1}} \setminus A_{k}(1)},
\end{equation}
where $A_k(1)$ is the trapezoid $A_k$ shifted by one unit to the right. 

We can repeat this process. 
Let us be more precise. 
If we set $\Upsilon_k(1)= \Upsilon_k-(k+2)\alpha_{i+2,i+2+(h-2)}$ and $\Upsilon_k'(1)= \Upsilon_k'-(k+2)\alpha_{i+2,i+2+(h-2)}$ then the conditions in \eqref{eq: upsilon } and \eqref{eq: upsilon prime} are ``shifted by one unit to the right''. 
Concretely, we have
\begin{equation} \label{eq: upsilon A}
    \pint{\Upsilon_k(1)}{\alpha_t} = \left\{  
    \begin{array}{rl}
      0,    & \mbox{if } i+1-k \leq t \leq i+1; \\
      -(k+2),    & \mbox{if } t=i+2; \\
      0,  & \mbox{if }  i+3\leq  t \leq j +\varkappa \mbox{ and } t\neq j+2;\\
      k+2, & \mbox{if }  t=j+2 \leq j+\varkappa.
    \end{array} \right.
\end{equation}
and 
\begin{equation} \label{eq: upsilon prime A}
    \pint{\Upsilon_k'(1)}{\alpha_t} = \left\{  
    \begin{array}{rl}
      0,    & \mbox{if } i+1-k \leq t \leq i+1; \\
      -(k+1),    & \mbox{if } t=i+2; \\
      0,  & \mbox{if }  i+3\leq  t \leq j +\varkappa \mbox{ and } t\neq j+2;\\
      k+1, & \mbox{if }  t=j+2 \leq j+\varkappa.
    \end{array} \right.
\end{equation}
So that we can apply \Cref{coro: push pyramid to the right} again. 
We get
\begin{equation}
\begin{array}{l}
     \bfM_{\Upsilon_k(1)}^{\Phi^{\succ \alpha_{a-h,a-1}} \setminus A_{k}(1)} = q^{k+2} \bfM_{\Upsilon_k(1)-(k+2)\alpha_{i+3,i+3+(h-2)}}^{\Phi^{\succ \alpha_{a-h,a-1}} \setminus A_{k}(2)},   \\[10pt]
  \bfM_{\Upsilon_k'(1)}^{\Phi^{\succ \alpha_{a-h,a-1}} \setminus A_{k}(1)} =q^{k+1} \bfM_{\Upsilon_k'(1)-(k+1)\alpha_{i+3,i+3+(h-2)}}^{\Phi^{\succ \alpha_{a-h,a-1}} \setminus A_{k}(2)},
\end{array}
\end{equation}
where $A_k(2)$ is the trapezoid $A_k$ shifted by two units to the right. 

By continuing this way, we eventually obtain 
\begin{equation}
    \begin{array}{l}
         \bfM_{\mu}^Y =(-1)^{k+1}q^{\mathfrak{R}_k}q^{(k+2)(h-1+\varkappa)} \bfM_{\Upsilon_k-\zeta_k}^{\Phi^{ \succ \alpha_{a-h,a-1}  } \setminus A_k(h-1+\varkappa) }, \\[10pt]
         \bfM_{\mu'}^Y =(-1)^{k+1}q^{\mathfrak{R}_k} q^{(k+1)(h-1+\varkappa)}\bfM_{\Upsilon_k'-\zeta_k'}^{\Phi^{ \succ \alpha_{a-h,a-1}  } \setminus A_k(h-1+\varkappa) },
    \end{array}
\end{equation}
where $A_k(h-1+\varkappa)$ is the trapezoid $A_k$ shifted by $h-1+\varkappa$ units to the right,  
\begin{equation}
    \zeta_k=  (k+2) \sum_{u=1}^{h-1+\varkappa}  \alpha_{i+1+u,i+1+u+(h-2)}  \qquad \mbox{and} \qquad   \zeta_k'=  (k+1) \sum_{u=1}^{h-1+\varkappa}  \alpha_{i+1+u,i+1+u+(h-2)} .
\end{equation}
Equation  \eqref{eq: relation P con v negrita} yields
\begin{equation}
\begin{array}{rl}
   ( \Upsilon_k-\zeta_k) -(\Upsilon_k'-\zeta_k')   & \displaystyle = \mu -\mu' - \sum_{u=1}^{h-1+\varkappa}  \alpha_{i+1+u,i+1+u+(h-2)} \\[10pt]
     &\displaystyle =  P_{j,h}^{h-1}(\lambda) -\mathbf{v}_{j,h}(\lambda) - \sum_{u=1}^{h-1+\varkappa}  \alpha_{i+1+u,i+1+u+(h-2)} \\[10pt]
     & =0.
\end{array}
\end{equation}
Thus, $\bfM_{\mu}^Y= q^{h-1+\varkappa} \bfM_{\mu'}^Y$.
Finally, by the  definition of $Y=\Phi^{\succ \alpha_{i,j}} \cup \{  \alpha_{a-h,a-1} \} $ and the equality $\alpha_{a-h,i}=\alpha_{a-h,a-1}+\alpha_{a,i}  $, we get
\begin{equation}
       \bfM_{P_{j,h}^{h-1}(\lambda) - \alpha_{a,i}}^{\succ \alpha_{i,j}} - q\bfM_{P_{j,h}^{h-1}(\lambda)-\alpha_{a-h,i}}^{\succ \alpha_{i,j}} =
       q^{h-1 + \varkappa } \left( \bfM_{\mathbf{v}_{j,h}(\lambda)-\alpha_{a,i}}^{\succ \alpha_{i,j}}  -q \bfM_{\mathbf{v}_{j,h}(\lambda)- \alpha_{a-h,i}}^{\succ \alpha_{i,j}} \right),
\end{equation}
which is equivalent to our claim.
\end{proof}

We recall that
\begin{equation}
    \varrho = \varrho_{i,h}(P_{j,h}^{h-1} (\lambda)) = \varrho_{i,h} (  \mathbf{v}_{j,h}(\lambda) ) = i -dh+1.
\end{equation}
Therefore, a repeated application of \Cref{claim complicado} starting with $a=i-h+1$ yields
\begin{equation} \label{eq: consequence claim}
\begin{array}{l}
  \qquad\quad\quad \bfM_{P_{j,h}^{h-1}(\lambda) - \alpha_{i-h+1,i}}^{\succ \alpha_{i,j}} - q^{\varkappa+h-1} \bfM_{\mathbf{v}_{j,h}(\lambda)-\alpha_{i-h+1,i}}^{\succ \alpha_{i,j}}  
     \\ = q^{d-1}  \left( \bfM_{P_{j,h}^{h-1}(\lambda)-\alpha_{i-dh+1,i}}^{\succ \alpha_{i,j}} - q^{\varkappa+h-1} \bfM_{\mathbf{v}_{j,h}(\lambda)- \alpha_{i-dh+1,i}}^{\succ \alpha_{i,j}} \right) \\
     = q^{d-1}  \left( \bfM_{P_{j,h}^{h}(\lambda)}^{\succ \alpha_{i,j}} - q^{\varkappa+h-1} \bfM_{P_{i,h}\mathbf{v}_{j,h}(\lambda)}^{\succ \alpha_{i,j}} \right).
\end{array}
\end{equation}  
Thus, \eqref{eq: ultimo P existe} follows by multiplying \eqref{eq: consequence claim} by $q^{d_1}$ and using the identities in \eqref{eq:degreesA} and \eqref{eq:degreesB}.
\end{proof}


\subsection{The second inverse decomposition}

We now turn to the proof of the \emph{Second Inverse Decomposition (SID)}, the central result of this section and the key to prove the positivity of the pre-canonical bases in the next section.

Before stating the result,  let us explain why the SID is needed in the present situation.
If we apply the First Inverse Decomposition in \Cref{prop: second version} to a weight $\lambda$, we can find that some of the weights indexing the $\bfM$-elements in the decomposition may fail to be dominant (disregarding the subtracted roots).
More precisely, the problematic terms are $P_{j,h}^{h-1}(\lambda)$ and $v_{j,h}^{h-1}(\lambda)$.
The SID addresses this issue by replacing these non-dominant weights with dominant ones.
This ensures that all weights involved remain dominant, which will allows us to use an inductive argument  in the proof of \Cref{Conj A}.

\begin{proposition}\label{prop: second version refinado}
Let $\lambda\in X^{+}$, $\alpha_{i,j}\in \Phi^{\geq 2}$ and $h=j-i+1$.  
Define $\beta_m=\alpha_{i-m,j-m}$ and 
\begin{equation}
    k=k_{j,h}(\lambda) \coloneq \min\{ 0 \leq s \leq h-2 \mid \pint{\lambda}{\alpha_{j-s}}>0\} .
\end{equation}

Then, we have
\begin{equation}\label{eq: SID ecuacion combi}
 \bfM_{\lambda}^{\succeq \alpha_{i,j}} = \left\{ \begin{array}{ll}
 \displaystyle    \bfM_{\lambda}^{\succ \alpha_{i,j}}
     - q^{d}\,\bfM_{\mathbf{P}_{j,h}(\lambda)}^{\succ \alpha_{i,j}}
     - \sum_{m=1}^{k}   q^{d_m} X_m , 
       & \mbox{if } k \mbox{ exists;}  \\
 \displaystyle   \bfM_{\lambda}^{\succ \alpha_{i,j}}
     - q^{d}\,\bfM_{\mathbf{P}_{j,h}(\lambda)}^{\succ \alpha_{i,j}}
     - \sum_{m=1}^{h-2}   q^{d_m} X_m
     - q^{e_1} \left(
        \bfM_{\mathbf{v}_{j,h}(\lambda)}^{\succ \alpha_{i,j}}
          - q^{e_2}\,  \bfM_{\mathbf{P}_{i,h}\mathbf{v}_{j,h}(\lambda)}^{\succ \alpha_{i,j}}
       \right) ,    & \mbox{otherwise;}
  \end{array}   \right.  
\end{equation}
where 
$d = D(\mathbf{P}_{j,h})(\lambda)$, 
$d_m = D(v_{j,h}^m)(\lambda)$, 
$e_1 = D(\mathbf{v}_{j,h})(\lambda)$, $ e_2 = D(\mathbf{P}_{i,h})(\mathbf{v}_{j,h}(\lambda))$ and 
\begin{equation} \label{eq: defi X in SID}
  X_m  =  \bfM_{v^{m}_{j,h}(\lambda)}^{\succ \alpha_{i,j}}
         -   q\,\bfM_{v^{m}_{j,h}(\lambda)-\beta_m}^{\succ \alpha_{i,j}}.   
\end{equation}
\end{proposition}



% \noindent
% If $k$ exists then we have
% \begin{equation}\label{eq: first refined version}
%     \!  \bfM_{\lambda}^{\succeq \alpha_{i,j}}\! \!  = \! \bfM_{\lambda}^{\succ \alpha_{i,j}} \!  - q^{d} \bfM_{\mathbf{P}_{j,h}(\lambda)}^{\succ \alpha_{i,j}}\! \!  -  \displaystyle\sum_{m=1}^{k} q^{d_m} \! \left( \bfM_{v^{m}_{j,h}(\lambda)}^{\succ \alpha_{i,j}}\!  -\!  q\bfM_{v^{m}_{j,h}(\lambda)-\beta_m}^{\succ \alpha_{i,j}}\right),
% \end{equation}
% where 
% $d = D(\mathbf{P}_{j,h})(\lambda)$, 
% $d_m = D(v_{j,h}^m)(\lambda)$.

% % \begin{equation}\label{eq: first refined version}
% %     \!  \bfM_{\lambda}^{\succeq \alpha_{i,j}}\! \!  = \! \bfM_{\lambda}^{\succ \alpha_{i,j}} \!  - q^{D(\mathbf{P}_{j,h})(\lambda)} \bfM_{\mathbf{P}_{j,h}(\lambda)}^{\succ \alpha_{i,j}}\! \!  -  \displaystyle\sum_{m=1}^{k} q^{D(v^{m}_{j,h})(\lambda)} \! \left( \bfM_{v^{m}_{j,h}(\lambda)}^{\succ \alpha_{i,j}}\!  -\!  q\bfM_{v^{m}_{j,h}(\lambda)-\beta_m}^{\succ \alpha_{i,j}}\right).
% % \end{equation}
% \medskip
% \noindent If $k$ does not exist (this is, if $\pint{\lambda}{\alpha_{j-s}}=0$ for all $0\leq s \leq h-2$)  then we have
%  \begin{equation}\label{eq: second version refinada}
% \resizebox{0.91\textwidth}{!}{$
%  \displaystyle    \bfM_{\lambda}^{\succeq \alpha_{i,j}}
%     =  \bfM_{\lambda}^{\succ \alpha_{i,j}}
%      - q^{d}\,\bfM_{\mathbf{P}_{j,h}(\lambda)}^{\succ \alpha_{i,j}}
%      - \sum_{m=1}^{h-2}   q^{d_m} \left(
%         \bfM_{v^{m}_{j,h}(\lambda)}^{\succ \alpha_{i,j}}
%          -   q\,\bfM_{v^{m}_{j,h}(\lambda)-\beta_m}^{\succ \alpha_{i,j}}
%       \right)
%      - q^{e_1} \left(
%         \bfM_{\mathbf{v}_{j,h}(\lambda)}^{\succ \alpha_{i,j}}
%           - q^{e_2}\,  \bfM_{\mathbf{P}_{i,h}\mathbf{v}_{j,h}(\lambda)}^{\succ \alpha_{i,j}}
%        \right),
% $}
% \end{equation}
% where 
% $d = D(\mathbf{P}_{j,h})(\lambda)$, 
% $d_m = D(v_{j,h}^m)(\lambda)$, 
% $e_1 = D(\mathbf{v}_{j,h})(\lambda)$ and $ e_2 = D(\mathbf{P}_{i,h})(\mathbf{v}_{j,h}(\lambda))$.     


\begin{remark}\rm\label{remark operator not defined refinada}  
      We recall our convention that if some $\mathbf{P}$-operator, $v^m$-operator or $\mathbf{v}$-operator is not defined, then the corresponding $\bfM$-element is set equal to zero. 
\end{remark}

\begin{proof}
We divide the proof into two cases according to whether $k$ exists.

\textbf{Case I. } The integer $k$ exists. 

Applying \Cref{prop: second version} to $\lambda$, $\alpha_{i,j}$ and $k$, we obtain:
\begin{equation}\label{eq: demo first refined version 1}
    \!  \bfM_{\lambda}^{\succeq \alpha_{i,j}}\! \!  = \! \bfM_{\lambda}^{\succ \alpha_{i,j}} \!  - q^{D(P^{k}_{j,h})(\lambda)+1} \bfM_{P^{k}_{j,h}(\lambda) - \beta_{k}}^{\succ \alpha_{i,j}}\! \!  -  \displaystyle\sum_{m=1}^{k} q^{D(v^{m}_{j,h})(\lambda)} \! \left( \bfM_{v^{m}_{j,h}(\lambda)}^{\succ \alpha_{i,j}}\!  -\!  q\bfM_{v^{m}_{j,h}(\lambda)-\beta_m}^{\succ \alpha_{i,j}}\right).
\end{equation}

Comparing this with the target equality  in \eqref{eq: SID ecuacion combi}, it is enough to show the following identity:
\begin{equation}\label{eq: demo first refined version 2}
    \bfM_{P^{k}_{j,h}(\lambda) - \beta_{k}}^{\succ \alpha_{i,j}}  =  \begin{cases}
        q^{D(P_{j-k,h})(P^{k}_{j,h}(\lambda))-1}\bfM_{\mathbf{P}_{j,h}(\lambda)}^{\succ \alpha_{i,j}},& \text{if }\mathbf{P}_{j,h}(\lambda) \text{ is defined};\\
        0,&\text{otherwise}.
    \end{cases}
\end{equation}

We proceed by analyzing these two subcases.

\medskip
\noindent
\textbf{Case IA.} The weight $\mathbf{P}_{j,h}(\lambda)$ is defined.

We first notice that 
\begin{equation}\label{eq: Pk+1 = P negrita}
    \pint{P^{k+1}_{j,h}(\lambda)}{\alpha_{j-k}} = \pint{\lambda}{\alpha_{j-k}} - \sum_{m=j-k}^{j}\pint{\alpha_{\varrho_{m},m}}{\alpha_{j-k}} = \pint{\lambda}{\alpha_{j-k}} - 1 \geq 0,
\end{equation}
where $\varrho_{m} = \varrho_{m,h}(P^{j-m}_{j,h}(\lambda))$. The final equality follows from \Cref{remark producto interno} and the definition of the integer $k = k_{j,h}(\lambda)$. 
From \eqref{eq: Pk+1 = P negrita}, we conclude that $\mathbf{P}_{j,h}(\lambda) = P^{k+1}_{j,h}(\lambda)$.


If $\pint{\lambda}{\alpha_{i-k}} \geq 1$, then by the definition of the $P$-operator, $P^{k+1}_{j,h}(\lambda)= P^{k}_{j,h}(\lambda) - \beta_{k}$.
The result follows via substituting this directly into \eqref{eq: demo first refined version 2} and noticing that $D(P_{j-k,h})(P^{k}_{j,h}(\lambda)) = 1$.

We now prove the equality under the assumption $\pint{\lambda}{\alpha_{i-k}} = 0$. 
By applying \Cref{lem: ceros por Ps} to $\lambda$, $\alpha_{i,j}$, and $k+1$, we get
\begin{equation}\label{eq: demo first refined version 3}
    \pint{P^{k}_{j,h}(\lambda)}{\alpha_{t}} = 0 \quad \mbox{for all}\; t\in T=\bigcup_{m=0}^{d} [i-k-mh,i-mh],
\end{equation}
 where $d$ is the unique integer satisfying $\varrho_{j-k,h}(P^{k}_{j,h}(\lambda)) = i-k-(d+1)h$.
From this, it follows that
\begin{equation} \label{eq: demo first refined version 4}
    \pint{P^{k}_{j,h}(\lambda) - \beta_{k}}{\alpha_{t}}  = \begin{cases}
            0, & \mbox{if } t \in T \setminus \{i-k\};\\
            -1, & \mbox{if } t=i-k.
    \end{cases}
\end{equation}
We can now apply \Cref{lem: free P} to $K=\Phi^{\succ \alpha_{i,j}}$, $\alpha_{i,j}$, $p=p'=i-k$, the set $T$, and $P^{k}_{j,h}(\lambda)-\beta_{k}$.
We obtain
\begin{equation}\label{eq: demo first refined version 5}
    \bfM_{P^{k}_{j,h}(\lambda) - \beta_{k}}^{\succ \alpha_{i,j}} = q^{D(P_{i-k-1,h})(P^{k}_{j,h}(\lambda) - \beta_{k})}\bfM_{P_{i-k-1,h}(P^{k}_{j,h}(\lambda)-\beta_{k})}^{\succ \alpha_{i,j}}
    = q^{\D{P,j-k,h}(P^{k}_{j,h}(\lambda))-1}\bfM_{P^{k+1}_{j,h}(\lambda)}^{\succ \alpha_{i,j}},
\end{equation}
where the second equality holds because
\begin{equation}\label{eq: demo first refined version 6}
\begin{array}{lcl}
    P_{i-k-1,h}(P^{k}_{j,h}(\lambda)-\beta_{k}) &=& P_{j-k,h}(P^{k}_{j,h}(\lambda)) = P^{k+1}_{j,h}(\lambda) \\[10pt]
    D(P_{i-k-1,h})(P^{k}_{j,h}(\lambda) - \beta_{k})+1 &=& \D{P,j-k,h}(P^{k}_{j,h}(\lambda)).
\end{array}
\end{equation}
This establishes the first case of \eqref{eq: demo first refined version 2}.

\medskip
\noindent
\textbf{Case IB.  } The weight $\mathbf{P}_{j,h}(\lambda)$ is not defined.

We must show that $\bfM_{P^{k}_{j,h}(\lambda) - \beta_{k}}^{\succ \alpha_{i,j}} = 0$. 

First, if $P^{k}_{j,h}(\lambda)$ is not defined, then by convention  $\bfM_{P^{k}_{j,h}(\lambda) - \beta_{k}}$  vanishes and we are done.  

Second, assume $P^{k}_{j,h}(\lambda)$ is defined. 
Since $\mathbf{P}_{j,h}(\lambda)$ is not defined, the definition of $k$ forces  
that $P^{k+1}_{j,h}(\lambda)$ is not defined.
Applying \Cref{lem: ceros por Ps} to $\lambda$, $\alpha_{i,j}$, and $k$ gives $\pint{P^{k-1}_{j,h}(\lambda)}{\alpha_{t}} = 0$ for all $t\in T$, where
\begin{equation}
    T=\bigcup_{m=0}^{d}[i-k+1-mh,i-mh],
\end{equation}
and $d$ is the unique integer satisfying $\varrho_{k} = \varrho_{j-k+1,h}(P^{k-1}_{j,h}(\lambda)) = i-k+1-(d+1)h$.
Since $P^{k+1}_{j,h}(\lambda)$ is not defined, \Cref{lem: desigualdad de varrhos} implies $\varrho_{k} = 1$. In other words, $P^{k}_{j,h}(\lambda) = P^{k-1}_{j,h}(\lambda) -\alpha_{1,j-k+1}$.
By \Cref{remark producto interno}, we conclude that $\pint{P^{k}_{j,h}(\lambda)}{\alpha_{t}}=0$ for all $t\in T$.

Furthermore, the condition that $P^{k+1}_{j,h}(\lambda)$ is not defined also implies, by the definition of the $P$-operator, that $\pint{P^{k}_{j,h}(\lambda)}{\alpha_{t}} =0$ for all $t\in \{i-k-mh\mid 0\leq m\leq d\}$.
Combining these zero-conditions, we conclude that $\pint{P^{k}_{j,h}(\lambda)}{\alpha_{t}} =0$ for all $t\in T'$, where
\begin{equation}
    T'=\bigcup_{m=0}^{d} [i-k-mh,i-mh].
\end{equation}
Using \Cref{remark producto interno} again, we obtain
\begin{equation}
    \pint{P^{k}_{j,h}(\lambda) - \beta_{k}}{\alpha_{t}} = \begin{cases}
            0,&\text{if } t\in T'\setminus \{i-k\};\\
            -1,&\text{if } t = i-k.
    \end{cases}
\end{equation}
As in the first case, we apply \Cref{lem: free P} to $K=\Phi^{\succ \alpha_{i,j}}$, $\alpha_{i,j}$, $p=p'=i-k$, the set $T'$, and $P^{k}_{j,h}(\lambda)-\beta_{k}$, to obtain $\bfM_{P^{k}_{j,h}(\lambda)-\beta_{k}}^{\succ \alpha_{i,j}}=0$, as we wanted to show. 

This completes the proof of \textbf{Case I}.


\bigskip
\textbf{Case II. } The integer $k$ does not exist. 

We notice that the hypothesis in this case implies that   $\pint{\lambda}{\alpha_{j-s}}=0$ for all $0\leq s \leq h-2$.

We only prove the case when all the weights involved in \eqref{eq: SID ecuacion combi} exist, the other cases follow in a similar fashion. 

By \Cref{prop: second version} applied to $ h-1$ we have
\begin{equation} \label{eq:loquedaFID}
     \bfM_{\lambda}^{\succeq \alpha_{i,j}} = \bfM_{\lambda}^{\succ \alpha_{i,j}} - q^{D(P^{h-1}_{j,h})(\lambda)+1} \bfM_{P^{h-1}_{j,h}(\lambda) - \beta_{h-1}}^{\succ \alpha_{i,j}} - \displaystyle\sum_{m=1}^{h-1} q^{d_m} \left( \bfM_{v^{m}_{j,h}(\lambda)}^{\succ \alpha_{i,j}} - q\bfM_{v^{m}_{j,h}(\lambda)-\beta_m}^{\succ \alpha_{i,j}}\right).
\end{equation}
By comparing \eqref{eq: SID ecuacion combi} and \eqref{eq:loquedaFID} it is easy to note that both formulas differ by three terms. 
We work with these terms. 
We have

\begin{equation}\label{eq: pasar de FID to SID}
\begin{array}{l}
      \quad   q^{D(P^{h-1}_{j,h})(\lambda)+1} \bfM_{P^{h-1}_{j,h}(\lambda) - \beta_{h-1}}^{\succ \alpha_{i,j}}
     + q^{d_{h-1}} \left( \bfM_{v^{m}_{j,h}(\lambda)}^{\succ \alpha_{i,j}} - q\bfM_{v^{m}_{j,h}(\lambda)-\beta_{h-1}}^{\succ \alpha_{i,j}} \right)    \\[15pt]
     = q^{D(P^{h-1}_{j,h})(\lambda)+1} \bfM_{P^{h-1}_{j,h}(\lambda) - \beta_{h-1}}^{\succ \alpha_{i,j}}
     + q^{e_1} \left( \bfM_{\mathbf{v}_{j,h}(\lambda)}^{\succ \alpha_{i,j}} - q\bfM_{\mathbf{v}_{j,h}(\lambda)-\beta_{h-1}}^{\succ \alpha_{i,j}} \right)   \\[15pt]
     = q^{D(P^{h}_{j,h})(\lambda)} \bfM_{P^{h}_{j,h}(\lambda)}^{\succ \alpha_{i,j}}
     + q^{e_1} \bfM_{\mathbf{v}_{j,h}(\lambda)}^{\succ \alpha_{i,j}} - q^{D(P_{i,h}\mathbf{v}_{j,h})(\lambda)}\bfM_{P_{i,h}\mathbf{v}_{j,h}(\lambda)}^{\succ \alpha_{i,j}}     \\[15pt]
     = q^{d} \bfM_{\mathbf{P}_{j,h}(\lambda)}^{\succ \alpha_{i,j}}
     + q^{e_1} \bfM_{\mathbf{v}_{j,h}(\lambda)}^{\succ \alpha_{i,j}} - q^{e_1+e_2}\bfM_{\mathbf{P}_{i,h}\mathbf{v}_{j,h}(\lambda)}^{\succ \alpha_{i,j}}  . 
\end{array}
\end{equation}


Let us briefly indicate the origin of these equalities.
The first equality is a consequence of \Cref{lem: completar Q en caso h-1}.
The second equality follows from \Cref{lem: ultimo P}.
The last equality is implied by \Cref{lem: bfP}.
Finally, let us note that the equalities appearing in the exponents come directly from the definition of the operators involved and the additivity in the  corresponding degrees.

Then, \eqref{eq: SID ecuacion combi} is obtained by replacing \eqref{eq: pasar de FID to SID} in \eqref{eq:loquedaFID}.
\end{proof}


% \newpage


% \begin{proposition}\label{prop: second version refinado}   
%     Let $\lambda\in X^{+}$, $\alpha_{i,j}\in \Phi^{\geq 2}$ and $h=j-i+1$. 
%     Suppose that $\langle \lambda , \alpha_r \rangle  =0$, for all $i+1\leq r \leq j$. 
%     Furthermore, write  $\beta_r=\alpha_{i-r,j-r}$. 
%     Then, we have    
%     \begin{equation}\label{eq: second version refinada}
% \resizebox{0.91\textwidth}{!}{$
%  \displaystyle    \bfM_{\lambda}^{\succeq \alpha_{i,j}}
%     =  \bfM_{\lambda}^{\succ \alpha_{i,j}}
%      - q^{d}\,\bfM_{\mathbf{P}_{j,h}(\lambda)}^{\succ \alpha_{i,j}}
%      - \sum_{r=1}^{h-2}   q^{d_r} \left(
%         \bfM_{v^{r}_{j,h}(\lambda)}^{\succ \alpha_{i,j}}
%          -   q\,\bfM_{v^{r}_{j,h}(\lambda)-\beta_r}^{\succ \alpha_{i,j}}
%       \right)
%      - q^{e_1} \left(
%         \bfM_{\mathbf{v}_{j,h}(\lambda)}^{\succ \alpha_{i,j}}
%           - q^{e_2}\,  \bfM_{\mathbf{P}_{i,h}\mathbf{v}_{j,h}(\lambda)}^{\succ \alpha_{i,j}}
%        \right),
% $}
% \end{equation}
% where 
% $d = D(\mathbf{P}_{j,h})(\lambda)$, 
% $d_r = D(v_{j,h}^r)(\lambda)$, 
% $e_1 = D(\mathbf{v}_{j,h})(\lambda)$ and $ e_2 = D(\mathbf{P}_{i,h})(\mathbf{v}_{j,h}(\lambda))$. 
% \end{proposition}


% \begin{remark}\rm\label{remark operator not defined refinada}
%        We emphasize that if any $\mathbf{P}$-operator, $v^r$-operator or $\mathbf{v}$-operator is not defined, the corresponding $\bfM$-element is set equal to zero. 
% \end{remark}


% \begin{proof}
% We only prove the case when all the weights involved in \eqref{eq: second version refinada} exist, the other cases follow in a similar fashion. 

% By \Cref{prop: second version} applied to $k= h-1$ we have
% \begin{equation} \label{eq:loquedaFID}
%      \bfM_{\lambda}^{\succeq \alpha_{i,j}} = \bfM_{\lambda}^{\succ \alpha_{i,j}} - q^{D(P^{h-1}_{j,h})(\lambda)+1} \bfM_{P^{h-1}_{j,h}(\lambda) - \beta_{h-1}}^{\succ \alpha_{i,j}} - \displaystyle\sum_{r=1}^{h-1} q^{d_r} \left( \bfM_{v^{r}_{j,h}(\lambda)}^{\succ \alpha_{i,j}} - q\bfM_{v^{r}_{j,h}(\lambda)-\beta_r}^{\succ \alpha_{i,j}}\right).
% \end{equation}
% By comparing \eqref{eq: second version refinada} and \eqref{eq:loquedaFID} it is easy to note that both formulas differ by three terms. 
% We work with these terms. 
% We have

% \begin{equation}\label{eq: pasar de FID to SID}
% \begin{array}{l}
%       \quad   q^{D(P^{h-1}_{j,h})(\lambda)+1} \bfM_{P^{h-1}_{j,h}(\lambda) - \beta_{h-1}}^{\succ \alpha_{i,j}}
%      + q^{d_{h-1}} \left( \bfM_{v^{r}_{j,h}(\lambda)}^{\succ \alpha_{i,j}} - q\bfM_{v^{r}_{j,h}(\lambda)-\beta_{h-1}}^{\succ \alpha_{i,j}} \right)    \\[15pt]
%      = q^{D(P^{h-1}_{j,h})(\lambda)+1} \bfM_{P^{h-1}_{j,h}(\lambda) - \beta_{h-1}}^{\succ \alpha_{i,j}}
%      + q^{e_1} \left( \bfM_{\mathbf{v}_{j,h}(\lambda)}^{\succ \alpha_{i,j}} - q\bfM_{\mathbf{v}_{j,h}(\lambda)-\beta_{h-1}}^{\succ \alpha_{i,j}} \right)   \\[15pt]
%      = q^{D(P^{h}_{j,h})(\lambda)} \bfM_{P^{h}_{j,h}(\lambda)}^{\succ \alpha_{i,j}}
%      + q^{e_1} \bfM_{\mathbf{v}_{j,h}(\lambda)}^{\succ \alpha_{i,j}} - q^{D(P_{i,h}\mathbf{v}_{j,h})(\lambda)}\bfM_{P_{i,h}\mathbf{v}_{j,h}(\lambda)}^{\succ \alpha_{i,j}}     \\[15pt]
%      = q^{d} \bfM_{\mathbf{P}_{j,h}(\lambda)}^{\succ \alpha_{i,j}}
%      + q^{e_1} \bfM_{\mathbf{v}_{j,h}(\lambda)}^{\succ \alpha_{i,j}} - q^{e_1+e_2}\bfM_{\mathbf{P}_{i,h}\mathbf{v}_{j,h}(\lambda)}^{\succ \alpha_{i,j}}  . 
% \end{array}
% \end{equation}


% Let us briefly indicate the origin of these equalities.
% The first equality is a consequence of \Cref{lem: completar Q en caso h-1}.
% The second equality follows from \Cref{lem: ultimo P}.
% The last equality is implied by \Cref{lem: bfP}.
% Finally, let us note that the equalities appearing in the exponents come directly from the definition of the operators involved and their corresponding degrees.

% Then, \eqref{eq: second version refinada} is obtained by replacing \eqref{eq: pasar de FID to SID} in \eqref{eq:loquedaFID}.

% \end{proof}




% % We now state a refined version of \Cref{prop: second version} that isolates the dominant term $P^k_{j,h}(\lambda)$. This term, which will later be denoted by $\mathbf{P}_{j,h}(\lambda)$ (cf. \Cref{def: P negrita}), is essential for proving several cases of the positivity conjecture in \Cref{section: positivity}.

% \begin{proposition}\label{prop: first refined version}
% \normalfont % Use \normalfont instead of \rm
% Let $\lambda\in X^{+}$, $\alpha_{i,j}\in \Phi^{\geq 2}$, and $h=j-i+1$. Let $k$ be an integer with $1\leq k <h$ satisfying $\pint{\lambda}{\alpha_r} = 0$ for all $j-k\leq r \leq j$, and $\pint{\lambda}{\alpha_{j-k+1}}\geq 1$.
% Set $\beta_m = \alpha_{i-m,j-m}$.
% Then, we have
% \begin{equation}\label{eq: first refined version}
%     \!  \bfM_{\lambda}^{\succeq \alpha_{i,j}}\! \!  = \! \bfM_{\lambda}^{\succ \alpha_{i,j}} \!  - q^{D(P^{k}_{j,h})(\lambda)+1} \bfM_{P^k_{j,h}(\lambda)}^{\succ \alpha_{i,j}}\! \!  -  \displaystyle\sum_{m=1}^{k-1} q^{D(v^{m}_{j,h})(\lambda)} \! \left( \bfM_{v^{m}_{j,h}(\lambda)}^{\succ \alpha_{i,j}}\!  -\!  q\bfM_{v^{m}_{j,h}(\lambda)-\beta_m}^{\succ \alpha_{i,j}}\right).
% \end{equation}
% \end{proposition}

% \begin{proof}
% Applying \Cref{prop: second version} to $\lambda$, $\alpha_{i,j}$ and $k-1$, we obtain:
% \begin{equation}\label{eq: demo first refined version 1}
%     \!  \bfM_{\lambda}^{\succeq \alpha_{i,j}}\! \!  = \! \bfM_{\lambda}^{\succ \alpha_{i,j}} \!  - q^{D(P^{k-1}_{j,h})(\lambda)+1} \bfM_{P^{k-1}_{j,h}(\lambda) - \beta_{k-1}}^{\succ \alpha_{i,j}}\! \!  -  \displaystyle\sum_{m=1}^{k-1} q^{D(v^{m}_{j,h})(\lambda)} \! \left( \bfM_{v^{m}_{j,h}(\lambda)}^{\succ \alpha_{i,j}}\!  -\!  q\bfM_{v^{m}_{j,h}(\lambda)-\beta_m}^{\succ \alpha_{i,j}}\right).
% \end{equation}
% Comparing this with the target equality \eqref{eq: first refined version}, the proposition will be proven if we can establish the following identity:
% \begin{equation}\label{eq: demo first refined version 2}
%     \bfM_{P^{k-1}_{j,h}(\lambda) - \beta_{k-1}}^{\succ \alpha_{i,j}}  =  \begin{cases}
%         q^{D(P_{j-k+1,h})(P^{k-1}_{j,h}(\lambda))-1}\bfM_{P^{k}_{j,h}(\lambda)}^{\succ \alpha_{i,j}},& \text{if }P^{k}_{j,h}(\lambda) \text{ is defined};\\
%         0,&\text{otherwise}.
%     \end{cases}
% \end{equation}
% We proceed by analyzing these two cases.

% \medskip
% \noindent
% \textbf{Case $P^k_{j,h}(\lambda)$ is defined.}
% We must establish the top line of \eqref{eq: demo first refined version 2}.

% If $\pint{\lambda}{\alpha_{i-k+1}} \geq 1$, then by the definition of the $P$-operator, $P^{k}_{j,h}(\lambda)= P^{k-1}_{j,h}(\lambda) - \beta_{k-1}$.
% The result follows via substituting this directly into \eqref{eq: demo first refined version 2} and noticing that $D(P_{j-k+1,h})(P^{k-1}_{j,h}(\lambda)) = 1$.

% We now prove the equality under the assumption $\pint{\lambda}{\alpha_{i-k+1}} = 0$. 
% By applying \Cref{lem: ceros por Ps} to $\lambda$, $\alpha_{i,j}$, and $k$, we get
% \begin{equation}\label{eq: demo first refined version 3}
%     \pint{P^{k-1}_{j,h}(\lambda)}{\alpha_{t}} = 0 \quad \mbox{for all}\; t\in T=\bigcup_{m=0}^{d} [i-k+1-mh,i-mh],
% \end{equation}
%  where $d$ is the unique integer satisfying $\varrho_{j-k+1,h}(P^{k-1}_{j,h}(\lambda)) = i-k+1-(d+1)h$.
% From this, it follows that
% \begin{equation} \label{eq: demo first refined version 4}
%     \pint{P^{k-1}_{j,h}(\lambda) - \beta_{k-1}}{\alpha_{t}}  = \begin{cases}
%             0, & \mbox{if } t \in T \setminus \{i-k+1\};\\
%             -1, & \mbox{if } t=i-k+1.
%     \end{cases}
% \end{equation}
% We can now apply \Cref{lem: free P} to $K=\Phi^{\succ \alpha_{i,j}}$, $\alpha_{i,j}$, $p=p'=i-k+1$, the set $T$, and $P^{k-1}_{j,h}(\lambda)-\beta_{k-1}$. This yields:
% \begin{equation}\label{eq: demo first refined version 5}
%     q\bfM_{P^{k-1}_{j,h}(\lambda) - \beta_{k-1}}^{\succ \alpha_{i,j}} = q^{\D{P,i-k,h}(P^{k-1}_{j,h}(\lambda) - \beta_{k-1})+1}\bfM_{P_{i-k,h}(P^{k-1}_{j,h}(\lambda)-\beta_{k-1})}^{\succ \alpha_{i,j}}
%     = q^{\D{P,j-k+1,h}(P^{k-1}_{j,h}(\lambda))}\bfM_{P^{k}_{j,h}(\lambda)}^{\succ \alpha_{i,j}},
% \end{equation}
% where the second equality holds because
% \begin{equation}\label{eq: demo first refined version 6}
% \begin{array}{lcl}
%     P_{i-k,h}(P^{k-1}_{j,h}(\lambda)-\beta_{k-1}) &=& P_{j-k+1,h}(P^{k-1}_{j,h}(\lambda)) = P^{k}_{j,h}(\lambda) \\[10pt]
%     \D{P,i-k,h}(P^{k-1}_{j,h}(\lambda) - \beta_{k-1})+1 &=& \D{P,j-k+1,h}(P^{k-1}_{j,h}(\lambda)).
% \end{array}
% \end{equation}
% It establishes the first case of \eqref{eq: demo first refined version 2}.

% \medskip
% \noindent
% \textbf{Case $P^k_{j,h}(\lambda)$ is not defined.}
% We must show that $\bfM_{P^{k-1}_{j,h}(\lambda) - \beta_{k-1}}^{\succ \alpha_{i,j}} = 0$. We consider two subcases.

% First, if $P^{k-1}_{j,h}(\lambda)$ is also not defined, then by \Cref{remark operator not defined}, the term $\bfM_{P^{k-1}_{j,h}(\lambda) - \beta_{k-1}}$ in \eqref{eq: demo first refined version 1} is zero. This trivially satisfies the second case of \eqref{eq: demo first refined version 2}.

% Second, assume $P^{k-1}_{j,h}(\lambda)$ is defined, but $P^{k}_{j,h}(\lambda)$ is not.
% Applying \Cref{lem: ceros por Ps} to $\lambda$, $\alpha_{i,j}$, and $k-1$ gives $\pint{P^{k-2}_{j,h}(\lambda)}{\alpha_{t}} = 0$ for all $t\in T$, where
% \begin{equation}
%     T=\bigcup_{m=0}^{d}[i-k+2-mh,i-mh],
% \end{equation}
% and $d$ is the unique integer satisfying $\varrho_{k-1} = \varrho_{j-k+2,h}(P^{k-2}_{j,h}(\lambda)) = i-k+2-(d+1)h$.
% Since $P^{k}_{j,h}(\lambda)$ is not defined, \Cref{lem: desigualdad de varrhos} implies $\varrho_{k-1} = 1$. In other words, $P^{k-1}_{j,h}(\lambda) = P^{k-2}_{j,h}(\lambda) -\alpha_{1,j-k+2}$.
% By \Cref{remark producto interno}, we conclude that $\pint{P^{k-1}_{j,h}(\lambda)}{\alpha_{t}}=0$ for all $t\in T$.

% Furthermore, the condition that $P^{k}_{j,h}(\lambda)$ is not defined also implies, by the definition of the $P$-operator, that $\pint{P^{k-1}_{j,h}(\lambda)}{\alpha_{t}} =0$ for all $t\in \{i-k+1-mh\mid 0\leq m\leq d\}$.
% Combining these zero-conditions, we conclude that $\pint{P^{k-1}_{j,h}(\lambda)}{\alpha_{t}} =0$ for all $t\in T'$, where
% \begin{equation}
%     T'=\bigcup_{m=0}^{d} [i-k+1-mh,i-mh].
% \end{equation}
% Using \Cref{remark producto interno} again, we obtain
% \begin{equation}
%     \pint{P^{k-1}_{j,h}(\lambda) - \beta_{k-1}}{\alpha_{t}} = \begin{cases}
%             0,&\text{if } t\in T'\setminus \{i-k+1\};\\
%             -1,&\text{if } t = i-k+1.
%     \end{cases}
% \end{equation}
% As in the first case, we apply \Cref{lem: free P} to $K=\Phi^{\succ \alpha_{i,j}}$, $\alpha_{i,j}$, $p=p'=i-k+1$, the set $T'$, and $P^{k-1}_{j,h}(\lambda)-\beta_{k-1}$, to obtain $\bfM_{P^{k-1}_{j,h}(\lambda)-\beta_{k-1}}^{\succ \alpha_{i,j}}=0$.

% Both cases of \eqref{eq: demo first refined version 2} are now established, which completes the proof.
% \end{proof}